\documentclass{beamer}
% Computer Science 1 — shared Beamer preamble.
% \input this file after \documentclass{beamer} in every Monday lecture
% deck so all decks share one visual system. Do not fork per-week copies of
% this file — edit it once here and every deck picks up the change on next
% compile. Vendored from professional_minds/presentations/beamer/theme/
% preamble.tex (same pattern semester_kickoff_week's own vendored copy
% uses), kept local rather than \input across repos so this repo stays
% self-contained and usable without a professional_minds checkout.

\usetheme{Boadilla}
\usecolortheme{seahorse}
\setbeamertemplate{navigation symbols}{}
\setbeamertemplate{footline}[frame number]

% Speaker notes: \note{...} on any frame. Compiling with the default
% `pdflatex deck.tex` produces a slides-only PDF (notes are embedded but
% hidden). Compiling with
% `pdflatex "\PassOptionsToPackage{show}{pgfpages}\input{deck.tex}"`
% or uncommenting the line below produces an instructor copy with each
% slide's notes on a second page — see README.md for both commands.
% \setbeameroption{show notes on second screen=right}

\usepackage{pgfpages}

\newcommand{\SessionTitleSlide}[4]{%
  % #1 = Week/Day (e.g. "Week 2 -- Monday"), #2 = session title,
  % #3 = essential question, #4 = primary source citation
  \begin{frame}
    \titlepage
  \end{frame}
  \begin{frame}{Essential Question}
    \Large #3

    \vspace{1em}
    \normalsize\emph{#4}
  \end{frame}
}


\title{Files, Terminals, and Your First Python Program}
\subtitle{CS1 Week 3 — Monday, August 31, 2026}
\author{Computer Science I}
\date{}

% Screenshots are real, de-identified frames from the class recording.
\newcommand{\classshot}[2]{%
  \begin{frame}{From class: #2}
    \begin{center}
      \includegraphics[width=.98\linewidth,height=.82\textheight,keepaspectratio]{img/#1}
    \end{center}
  \end{frame}%
}

\begin{document}
\begin{frame}
  \titlepage
\end{frame}

\begin{frame}{The lesson in one sentence}
  \Large
  The file manager and terminal are two views of the same filesystem; the computer follows exact instructions.
  \vfill
  \normalsize Predict. Try one small change. Read the result. Verify.
\end{frame}

\begin{frame}{Two views, one location}
  \begin{columns}[T]
    \column{.48\textwidth}
    \textbf{File manager}
    \begin{itemize}
      \item Visual folders and files
      \item Create and open items
      \item Save changes explicitly
    \end{itemize}
    \column{.48\textwidth}
    \textbf{Terminal}
    \begin{itemize}
      \item Text commands
      \item Exact current location
      \item Ask it to list or inspect
    \end{itemize}
  \end{columns}
\end{frame}

\classshot{0540_ls_cd_git}{the same \texttt{git} folder, in the file manager and in PowerShell}

\begin{frame}[fragile]{Where am I? What is here?}
\begin{verbatim}
ls
cd GIT
\end{verbatim}
\vspace{1em}
\begin{itemize}
  \item \texttt{ls} lists the current directory.
  \item \texttt{cd} changes the current directory.
  \item After creating a file, list again: the terminal will not guess.
\end{itemize}
\end{frame}

\classshot{0780_save_txt_unsaved}{\texttt{Hello\_class.txt} now shows up in the folder after we saved it}

\begin{frame}{Files have jobs and names}
  \begin{center}
  \begin{tabular}{lll}
    \textbf{Extension} & \textbf{Kind} & \textbf{Use} \\
    \hline
    \texttt{.txt} & plain text & notes or simple text \\
    \texttt{.md} & Markdown & structured readable notes \\
    \texttt{.py} & Python source & instructions Python can execute
  \end{tabular}
  \end{center}
  \vspace{1em}
  A \texttt{.py} file is text on disk; it becomes a program when its contents and execution context are correct.
\end{frame}

\begin{frame}[fragile]{Your first tiny Python file}
\begin{verbatim}
# A short note for another human
print("hello world")
\end{verbatim}
\begin{itemize}
  \item \texttt{\#} begins a comment: useful to humans, ignored by Python.
  \item \texttt{print()} is a function that sends text to output.
\end{itemize}
\end{frame}

\classshot{1380_comments_hash}{a real comment line inside the file, read back with \texttt{cat}}

\classshot{2820_print_function}{naming the parts: the function, the argument, the string}

\begin{frame}{Literal machine, useful partner}
  \begin{itemize}
    \item Python does not infer your intention.
    \item Names, punctuation, capitalization, and location matter.
    \item Once you learn the language, the machine becomes a resource instead of a mystery.
  \end{itemize}
\end{frame}

\begin{frame}[fragile]{Case sensitivity is a rule}
\begin{verbatim}
print("hello world")   # works
Print("hello world")   # different name: error
\end{verbatim}
\vspace{1em}
Humans see the intended word. Python sees two different spellings.
\end{frame}

\classshot{1800_case_sensitive_print}{\texttt{Print} vs \texttt{print} --- ``\,name 'Print' is not defined. Did you mean: 'print'?\,''}

\begin{frame}{Errors are information}
  \begin{enumerate}
    \item Read the error and locate the line.
    \item Predict one small correction.
    \item Change one thing and run again.
    \item Ask a person or tool to translate unfamiliar language.
  \end{enumerate}
  \vfill
  An error costs a little time; the correction teaches you what the rule is.
\end{frame}

\classshot{2100_powershell_herestring}{one small fix, then run it: \texttt{python .\textbackslash hello\_in\_python.py}}

\begin{frame}{Small tools from the lecture}
  \begin{itemize}
    \item \textbf{Up arrow:} recall the previous terminal command.
    \item \textbf{Tab:} complete a name and check available choices.
    \item \textbf{Ctrl+S:} save the file before testing it.
    \item \textbf{File inspection:} look at the text before asking what it does.
  \end{itemize}
\end{frame}

\classshot{2970_syntax_error_fix}{experimenting: \texttt{TWO\_PRINTS.py} and the repeated output it makes}

\begin{frame}{Before Wednesday}
  \begin{itemize}
    \item Save a tiny \texttt{.py} file with one comment and one \texttt{print()} line.
    \item Predict its output before running it.
    \item Wednesday: paired reasoning with the Week 3 material.
    \item Friday: show-and-tell/check-in on what worked and what you learned.
  \end{itemize}
\end{frame}

\begin{frame}{Takeaway}
  \Huge
  \centering
  Be precise, be curious, and use errors as evidence.
\end{frame}
\end{document}
